% chapters/theory/09.tex

\chapter{Trends and Open Directions in Foundation Models}
\label{chap:trends-open-directions}

The core ideas in this book—tokenization, embeddings, self-attention, Transformers, scaling, alignment, reasoning, retrieval, tools, and evaluation—describe the mainstream of modern language AI. Yet the field is still moving quickly. Many of the largest improvements come not from a single “new model,” but from steady refinement across architecture, training, data, systems, and evaluation. This chapter closes the theory section by connecting the pieces into a coherent picture and then surveying several directions that have been particularly active in 2025: multimodal Transformers, diffusion-style text generation, improved optimizers and training recipes, the rising importance of data quality, efficiency and small models, hardware-aware design, and the gradual democratization of agentic workflows. The goal is not to predict the future with certainty, but to provide a principled map of what is changing and why.

\section{A unifying picture: what stays constant}
It is helpful to notice what has \emph{not} changed. The Transformer remains the backbone of most foundation models. Its essential mechanism is self-attention: build token representations by letting tokens attend to one another through learned query, key, and value projections. Once attention is in place, most of the engineering and research questions become questions of refinement: how to represent position, how to stabilize training, how to make attention efficient at long contexts, how to tune generation and alignment, and how to connect the model to external information and actions.

The training pipeline also remains structurally stable: large-scale pretraining with a next-token objective produces a base model; supervised fine-tuning produces a helpful assistant; preference tuning aligns that assistant with human preferences; reinforcement learning with verifiable rewards produces reasoning behavior; retrieval and tools compensate for static knowledge and enable action. These steps can be recombined in new ways, but the overall story persists.

\section{Transformers beyond text: vision Transformers and multimodality}
Transformers were introduced for translation, but self-attention is not intrinsically tied to language. Attention operates on vectors. If we can represent other modalities—images, audio, video—as sequences of vectors, we can apply the same architecture.

The vision Transformer (ViT) demonstrates this clearly. An image is divided into patches. Each patch is mapped to a vector (a learned projection of pixel values). Positional information is added to preserve spatial structure. The resulting patch tokens are processed by a Transformer encoder. A special classification token serves as a pooled representation, analogous to BERT’s \texttt{[CLS]} token. This representation is then projected to a class label. The surprising result of ViT is that with enough data, this low-inductive-bias architecture can compete with and even surpass convolutional networks that hard-code locality through receptive fields.

Multimodal LLMs extend this idea to systems that accept both images and text. A common strategy is to use an image encoder to produce a sequence of “image tokens” and then concatenate them with text tokens as input to a decoder-only language model that generates text. Another strategy uses cross-attention to let text tokens attend to image representations within the network. Both approaches exploit the same core property: once everything is turned into token-like vectors, attention can fuse modalities.

A deeper lesson is that architectural tricks migrate across modalities. Rotary position embeddings (RoPE), originally designed for relative positioning in 1D token sequences, can be adapted to 2D spatial coordinates for images, and further adapted to mixed streams of image and text tokens. This kind of cross-pollination is a defining trait of the current era: progress in one modality often becomes a reusable tool in another.

\section{Diffusion-style language models}
The standard LLM is autoregressive: it generates one token at a time from left to right. This creates a fundamental bottleneck at inference: generation is sequential and cannot be parallelized across output tokens. Training is parallelizable via masking, but inference is not. As LLMs move into settings where latency matters—interactive coding, agents, real-time assistants—this sequential bottleneck becomes increasingly important.

Diffusion models, which became dominant in image generation, offer a different paradigm. In the image domain, diffusion generates by starting from noise and iteratively denoising. This works well because noise is easy to sample and has convenient mathematical properties. It also yields a controllable trade-off between speed and quality: more denoising steps typically improve quality.

Adapting diffusion to text is challenging because text is discrete. The common conceptual mapping is that “noise” for images corresponds to “masking” for text. Instead of adding Gaussian noise, one progressively replaces tokens with a \texttt{[MASK]} symbol until the sequence is fully masked. The reverse process is then trained to unmask tokens, reconstructing a coherent output conditioned on a prompt. Models in this family are sometimes described as masked diffusion models (MDMs) or diffusion-based LLMs.

A helpful intuition is to think of writing not as strictly left-to-right, but as iterative refinement. When humans draft, they often write a rough structure first and then refine sections. Diffusion-style decoding resembles this coarse-to-fine process: rather than committing to the first token and proceeding, the model can revise multiple positions over a small number of refinement steps.

The primary promise is speed. If a diffusion model requires a fixed number of refinement steps, and that number is much smaller than the output length, it can generate long outputs with fewer forward passes than autoregressive decoding. Claims of order-of-magnitude speedups appear in experimental reports, especially for longer outputs. Another practical advantage is bidirectional context usage within the generated region, which can benefit “fill-in-the-middle” code tasks.

The open question is whether diffusion language models can match frontier autoregressive models across the full range of capabilities: reasoning, tool use, alignment, and safety. Much of the existing ecosystem—decoding strategies, chain-of-thought training, RL pipelines—was built around autoregressive generation. Adapting these tools to diffusion is an active area of research.

\section{Architectural refinement is still ongoing}
It is tempting to treat the Transformer as a finished design, but many subcomponents remain in flux. Normalization moved from post-norm to pre-norm, and LayerNorm variants such as RMSNorm became common. Attention variants such as grouped-query attention balance quality and KV-cache efficiency. Sparse attention patterns and long-context mechanisms continue to evolve. Even activation functions are not fixed: ReLU-like functions such as GELU became common, and new activation functions continue to be proposed.

One area of refinement that is often underestimated is the optimizer. Adam became the default in deep learning because it works reliably, but new papers continue to challenge whether it is optimal for very large-scale training. New optimizers and stability tricks emerge as training runs push into regimes where small differences translate into large savings in compute and time. The broader lesson is that even “simple” components—optimizer choice, normalization placement, numerical formats—can have outsized effects at scale.

\section{Data is becoming the bottleneck}
Early scaling benefited from a relatively simple assumption: the internet contained vast amounts of human-generated text. That assumption is eroding. As LLMs become widespread, large portions of the web are now LLM-generated. Training on synthetic text can change the effective data distribution and reduce diversity, leading to a phenomenon commonly described as model collapse: the model learns from a narrower, less varied distribution and degrades over generations of self-training.

This does not imply that the field is doomed; it implies that data curation becomes a central competence. Instead of scraping everything, teams increasingly curate. A common pattern is to insert an intermediate “mid-training” stage between pretraining and fine-tuning: train on a large but higher-quality corpus tailored to desired domains, still using a pretraining-style objective. Synthetic data can still be useful, but it must be filtered and combined with human data carefully. The theme is shifting from “more data” to “better data.”

\section{Efficiency and the rise of small models}
For several years, progress was largely measured by benchmark improvements at the frontier. As models approach strong performance across many benchmarks, a second frontier becomes more important: efficiency. Serving LLMs is costly; inference dominates operational expenditure in many products. This motivates research into small language models (SLMs), distillation, quantization, sparse architectures, and improved decoding and caching strategies.

The key idea is the Pareto frontier. For a fixed cost budget—latency, dollars, memory—what is the best quality attainable? Different models occupy different points on this trade-off surface. In many real applications, the frontier is more relevant than the absolute best model, because the absolute best model may be too expensive or too slow.

\section{Hardware-aware thinking and new accelerators}
Transformers were not designed around a specific chip, but their success is tightly coupled to GPUs, which excel at matrix multiplication. As a result, the entire software stack is built around optimizing matrix-heavy operations, memory traffic, and caching. FlashAttention is an emblematic example: it does not change the model’s math, but it reorganizes computation to reduce memory movement because memory bandwidth is the real bottleneck.

A natural next step is to ask whether hardware itself should evolve to match Transformer workloads. Research prototypes explore accelerator designs that implement key operations more directly, sometimes even exploiting analog physical properties to reduce energy or latency. Whether such approaches become mainstream is uncertain, but the direction is clear: if models remain dominant, hardware will adapt to them, and model design will increasingly incorporate hardware constraints.

\section{Agents and the democratization of automation}
Tool calling and agentic workflows expand LLMs from “talkers” to “doers,” but they also amplify reliability and safety challenges. Each additional step in an agent’s loop introduces a chance of failure, and the probability of end-to-end success can decline quickly with length of the plan. This is why agent evaluation emphasizes reliability (for example, metrics that require all runs to succeed rather than at least one).

Despite these challenges, the long-term direction is toward democratization: enabling non-experts to build automations through natural language rather than code. As agents become more robust, they will increasingly mediate interaction with software systems, websites, and operating systems. This will require new standards and protocols, new safety mechanisms, and new evaluation regimes. The trajectory suggests that agentic interaction may become as common as search, but it will likely be gated by trust, permissions, and reliability.

\section{Staying current: how to keep learning}
The field moves quickly, so the skill is not memorizing today’s best model but learning how to stay current. Reading papers is necessary but not sufficient; implementations matter. Many research results become meaningful only when you inspect code and reproduce experiments. Benchmarks, blogs, and community discussions can help, but they also require skepticism: numbers can be gamed, and Goodhart’s law applies. The healthiest habit is to triangulate: compare paper claims, code, and real-world behavior on tasks that resemble your own.

\section{Summary}
Transformers remain the central architecture, but important improvements continue to arrive through refinement of positional methods, normalization, attention variants, optimizers, and training recipes. Multimodality shows how attention generalizes beyond text, and diffusion-style approaches show how paradigms from image generation may reshape text generation by reducing inference-time sequential bottlenecks. Data quality is increasingly central as synthetic text proliferates, and efficiency is becoming a dominant research axis as deployment costs matter. Hardware and software co-evolve with model demands, and agentic systems push LLMs into automation while creating new reliability and safety requirements. The overall picture is that the foundations are stable, but the frontier is still expanding—and it is expanding in multiple dimensions at once.

\section{Exercises}
\paragraph{1. Cross-modal reuse.}
Explain why a Transformer encoder can process both text tokens and image patches. What additional information must be added for images that is not needed for plain text tokens?

\paragraph{2. Autoregressive bottleneck.}
Why is inference-time generation fundamentally sequential in an autoregressive model? Why is training-time computation more parallelizable than inference-time?

\paragraph{3. Diffusion intuition.}
Describe the forward and reverse processes in diffusion for images, then explain the analogous processes for text using masking. What is the key challenge introduced by discrete tokens?

\paragraph{4. Data collapse.}
What does “model collapse” mean in the context of training on model-generated text? Why might diversity decrease, and why does that matter?

\paragraph{5. Pareto frontier thinking.}
Give an example where a smaller model is preferable to a larger one. What axes besides accuracy determine the “best” model for a product?
